\section{Alcance y Limitaciones}

\subsection{Alcance}

El presente informe incluye la definición de las bases de diseño, el cálculo del caudal de ventilación, la estimación de la potencia del ventilador, el dimensionamiento del área de impulsión y de las rejillas de exfiltración, la sustentación normativa de las 12 renovaciones de aire por hora, y la generación de datos de contorno para un modelo CFD. Adicionalmente, se presentan recomendaciones para el control de presión diferencial, la filtración del aire de impulsión y la estanqueidad del recinto.

\subsection{Limitaciones}

El estudio se limita al diseño conceptual del sistema de ventilación y presurización. No se incluye el diseño detallado de ductos, ya que la estrategia seleccionada emplea impulsión directa. No se realiza la selección comercial de un ventilador específico, aunque se indican los parámetros mínimos requeridos. El modelo CFD se ejecutó en Autodesk CFD con las condiciones de contorno presentadas en la Tabla \ref{tab:cfd_boundaries}; los resultados se discuten en las Secciones \ref{sec:resultados} y \ref{sec:analisis}. Las propiedades del aire se consideran aproximadamente constantes a 25 $^{\circ}$C y 101.3 kPa.
